\specialchapter{Epilogue}{A grammar of moves: when the sea rises, and when it cannot}{29}
\begin{refsection}[chapters/18-epilogue.bib]

Seventeen cases have taught one skill from seventeen directions. This epilogue puts the moves side by side, looks at how they combine, and then turns to the questions that a casebook of successes tends to avoid: when the sea cannot rise, and when it should not.

\addsec{The moves side by side}

The table below compresses each case to its move, its trigger, and the place where the difficulty went. Read down the last column: in every case the difficulty was paid once, in building something, and the payment was then spread over many problems. That is the signature of a dissolution, as opposed to a renaming.

\begin{fullwidth}
\small
\renewcommand{\arraystretch}{1.22}
\begin{tabularx}{\linewidth}{@{}r>{\raggedright}p{0.17\linewidth}>{\raggedright}p{0.33\linewidth}X@{}}
\toprule
 & \textsc{move} & \textsc{trigger} & \textsc{where the difficulty went}\\
\midrule
1 & completion & limits escape the space; theorems need uniform convergence & constructing Lebesgue measure\\
2 & completion & optimizers are sequences, not objects & measure theory and compactness of graphon space\\
3 & axiomatization & paradoxes with several ``correct'' answers & measure-theoretic foundations and extension theorems\\
4 & axiomatization & several proofs share one exchange lemma & verifying the exchange axiom for each new structure\\
5 & duality & ``for every adversary there is a response'' & convexity: the separating hyperplane theorem; integrality\\
6 & existence to dynamics & an equilibrium is known to exist but not how to find it & the regret analysis of one potential function\\
7 & transfer to topology & a lower bound resists counting; a hidden antipodal symmetry & the Borsuk--Ulam theorem\\
8 & diagonalize the symmetry & a probability is a polynomial in values of a function on a group & setting up the Fourier dictionary; hypercontractivity\\
9 & linearization & a set must be large but counting gives weak bounds & finding a vector space of the right dimension\\
10 & geometrization & an inequality compares volumes or entropies of two objects & existence and regularity of optimal maps\\
11 & the right invariant & a union bound over infinitely many events & the Sauer--Shelah lemma and symmetrization\\
12 & the right invariant & an adversary chooses the data & Ramsey-type arguments behind the equivalences\\
13 & localize to a limit & a pathology at a single point breaks a ``theorem'' & contiguity and local asymptotic normality\\
14 & exploit a free symmetry & validity is needed without a model & the exchangeability assumption itself\\
15 & encoding & a randomized process might run for a long time & designing a lossless encoding of the run\\
16 & proofs to algorithms & identifiability is provable but no algorithm is known & finding low-degree certificates; running time\\
17 & unification & parallel definitions with $\forall\exists$ structure & tuning definitions so that parameters translate\\
\bottomrule
\end{tabularx}
\end{fullwidth}

Three regularities stand out. First, the triggers are \emph{local} symptoms that you can notice in the middle of a proof attempt: an extra hypothesis that is used only to swap two limits, a lemma that appears in three papers with three proofs, an equilibrium that exists but cannot be computed, a union bound that runs out of room. A good working habit is to treat each such symptom as a question about the setting rather than as a technical obstacle to be overpowered.

Second, the move is almost always a change of \emph{object}, not a change of technique. Lebesgue did not find a cleverer way to estimate Riemann sums; he changed what an integral is. Littlestone did not find a cleverer adversary; he defined a tree that records every adversary at once. The technique of the old proofs often survives, relocated, inside the definitions of the new setting.

Third, the new object is usually humble. A function on the unit square, a list whose order does not matter, a polynomial that vanishes on a set, a binary tree of labelled points, a noodle. Grothendieck's advice to stoop low is not a matter of taste; it reflects the fact that the right object is the one that keeps exactly the information the problem uses, and such objects tend to be simple.

\addsec{How moves combine}

No case in this book uses only one move, and the most powerful dissolutions chain several.

\paragraph{Completion, then compactness, then existence.} Completing a space is rarely the end. In the completed space one proves compactness, and compactness turns optimization problems into existence theorems: extremal graphons exist (\cref{ch:graphon}), optimal transport plans exist (\cref{ch:transport}), minimizers of the Dirichlet energy exist in a Sobolev space (an exercise of \cref{ch:lebesgue}). The price is effectiveness. Compactness arguments say that an optimizer exists without saying how good a finite approximation is, which is why the combinatorial regularity lemma, with its tower-type bounds, was not made obsolete by graph limits.

\paragraph{Duality, then dynamics.} Duality produces certificates; dynamics produce them algorithmically. Linear programming duality proves the minimax theorem in a few lines (\cref{ch:lpduality}); two players running no-regret algorithms against each other prove it again and compute the equilibrium (\cref{ch:regret}). The same chain appears in the sum-of-squares method, where a duality between proofs and pseudo-distributions becomes an algorithm once the proofs have low degree (\cref{ch:sos}).

\paragraph{Symmetry, in three costumes.} The symmetrization argument of VC theory, the Fourier expansion on the Boolean cube and the rank argument of conformal prediction all exploit a group acting on the data (\cref{ch:vc,ch:fourier,ch:conformal}). In the first the group swaps a sample with a ghost sample, in the second it acts on the cube by translations, and in the third it permutes exchangeable observations. Recognizing which group acts, and on what, is often the whole move.

\paragraph{Encoding, in three costumes.} Moser's entropy compression (\cref{ch:compression}), Trevisan's reconstruction argument for extractors (\cref{ch:pseudorandom}) and the sample-compression view of learning (touched on in \cref{ch:vc}) share one idea: if a bad event lets you describe something random in fewer bits than it contains, the bad event is rare. The encoding is the move; the counting that follows is routine.

\paragraph{Invariants that agree.} When the same invariant appears independently in several places, it is probably the right one. The Littlestone dimension characterizes online learnability, private learnability and, in model theory, stability (\cref{ch:littlestone}). Such coincidences are evidence, not proof, but they are the best evidence available that a setting is natural rather than convenient.

\addsec{When the sea cannot rise}

Some problems resist every setting tried so far, and for a few of them there are theorems explaining why. These \term{barrier} results are dissolutions of a second order: they do not solve the problem, but they rule out whole families of reframings, and knowing them can save years.\index{barrier}

The question whether $\mathrm{P}=\mathrm{NP}$ is surrounded by three such barriers. \textcite{BakerGillSolovay1975} constructed oracles relative to which $\mathrm{P}=\mathrm{NP}$ and others relative to which they differ, so no proof technique that is insensitive to oracles can settle the question. \textcite{RazborovRudich1997} showed that a broad class of ``natural'' lower-bound arguments for circuits would, if it succeeded, break pseudorandom functions believed to exist. \textcite{AaronsonWigderson2009} extended the first barrier to algebraic techniques of the kind that proved $\mathrm{IP}=\mathrm{PSPACE}$. Each barrier describes a sea that has already been tried and a shore it cannot reach.

Independence results are barriers of a different kind. G\"odel and Cohen showed that the continuum hypothesis can be neither proved nor refuted from the usual axioms of set theory \parencite{Godel1940,Cohen1963}, so no setting expressible in those axioms decides it. Learning theory has its own example: \textcite{BenDavidEtAl2019} exhibited a learning problem whose learnability is independent of the same axioms, which is one reason the clean equivalences of \cref{ch:vc} are a special feature of binary classification rather than a law of learning. Impossibility theorems inside a field play the same role at a smaller scale; the impossibility of distribution-free conditional coverage in \cref{ch:conformal} tells you that no reframing of conformal prediction will deliver it, so effort should go into assumptions instead.

\addsec{When not to raise the sea}

The rising sea is a method, not a virtue, and it has costs. Building a setting takes time, often years and sometimes a generation, and the setting has to be learned by everyone who uses it. Three kinds of problems show that the chisel remains indispensable.

Some problems have been settled only by direct attack, with no dissolution in sight. The four colour theorem was proved by \textcite{AppelHaken1977} with a large computer-assisted case analysis, and the proof was later formalized completely by \textcite{Gonthier2008}; no conceptual setting has yet made it short. The Kakeya conjecture in three dimensions, whose finite-field analogue dissolved in a page under the polynomial method (\cref{ch:dimension}), was proved in a 2025 preprint by \textcite{WangZahl2025}, through a long argument of a very different kind that works across many scales at once. The finite-field sea did not reach the Euclidean shore.

Some problems have been reframed beautifully without dissolving. Sidorenko's conjecture \parencite{Sidorenko1993} is a clean inequality between homomorphism densities, and the graphon language of \cref{ch:graphon} states it in one line; it resisted proof for three decades. A preprint of September 2026 claims a counterexample \parencite{OpenAI2026}, found by exhibiting a kernel with a density gap and sampling a finite graph from it, which is a construction rather than a dissolution; when this book was written, the claim had not been independently checked. The $3x+1$ problem has been placed in dynamical, probabilistic and logical settings \parencite{Lagarias1985}, each illuminating and none decisive. A reframing that makes a problem easier to state is not the same as one that makes it easier to solve.

And some proofs are already short. A clean direct argument does not improve by being embedded in a theory, and readers pay for every layer of abstraction they must climb. The prologue's warning bears repeating: never soak a nut that is already open.

\addsec{A practice routine}

The moves in this book were found by people who were looking at a specific problem, not by people looking for moves. The following routine is an attempt to make that looking deliberate. None of it is a guarantee.

\begin{enumerate}[label=\arabic*.,leftmargin=*]
\item \emph{Keep a record of pain.} When a proof needs a hypothesis only to make one step work, or when the same lemma appears in several of your papers with different proofs, write it down. The triggers in the table above are entries from such a record.
\item \emph{Ask what the proof actually uses.} List every property of the objects that the argument touches. If the list is short, the right setting may be the class of all objects with those properties, as with matroids (\cref{ch:matroid}).
\item \emph{Look for escaping limits, hidden parameters, adversaries and free symmetries.} These four questions find most of the moves of Parts I to III.
\item \emph{Test a candidate setting on known results first.} A setting that cannot re-prove three solved problems more easily than their original proofs will not solve the open one. This is the dissolution test, applied before the investment.
\item \emph{Ask where the difficulty went.} If you cannot say where it went, it is probably still there, hidden in a definition or a hypothesis that is as hard to verify as the original problem.
\item \emph{Look for the cheapest falsifier of every unification.} A claim that two things are secretly the same is a conjecture until tested; the fastest way to test it is to look for a property one has and the other lacks.
\end{enumerate}

\begin{move}{Choose the setting before the technique}
\movepart{In one sentence} Before attacking a problem, ask in which world it is easy, and whether that world is worth building.
\movepart{Triggers} Every trigger in this book: escaping limits, recurring lemmas, adversaries, hidden symmetries, parameters that a paradox forces into the open, and existence without construction.
\movepart{Recipe} Isolate what the proof uses; propose an object that keeps exactly that; test it on solved problems; measure what got shorter; ask where the difficulty went.
\movepart{Anti-pattern} Building a world for its own sake, and proving the old theorems again at greater length in it.
\end{move}

\begin{exercises}
\item \research Return to the theorem you chose in the last exercise of the prologue. Apply the routine above: list what its proof uses, propose a setting, and test the setting on two solved results in the same area. Write a one-paragraph dissolution test for your proposal, including the honest entry ``none'' where it applies.
\item \research Pick one row of the table at the start of this epilogue and find a problem from a different field with the same trigger. Carry out the move, or explain precisely why it fails.
\item \research Barrier results describe shores that a class of settings cannot reach. Choose an open problem in your field and look for its barriers: oracles, independence, impossibility theorems, or lower bounds against a class of methods. What do they tell you about the setting a solution would need?
\item \research Find a case of your own in which a proof became shorter after a change of setting but the setting turned out to cost more than it saved. What made it fail the test?
\end{exercises}

\notesandrefs

The relativization, natural proofs and algebrization barriers are due to \textcite{BakerGillSolovay1975}, \textcite{RazborovRudich1997} and \textcite{AaronsonWigderson2009}. G\"odel's and Cohen's work on the continuum hypothesis is in \textcite{Godel1940} and \textcite{Cohen1963}, and the undecidability of a learning problem in \textcite{BenDavidEtAl2019}. The four colour theorem is due to \textcite{AppelHaken1977}, with a formal proof by \textcite{Gonthier2008}; the three-dimensional Kakeya conjecture was proved by \textcite{WangZahl2025}; Sidorenko's conjecture is from \textcite{Sidorenko1993}, and the claimed counterexample from \textcite{OpenAI2026}; and \textcite{Lagarias1985} surveys the $3x+1$ problem. The image of the rising sea, and its limits as a method, are discussed by \textcite{McLarty2007}; P\'olya's heuristics \parencite{Polya1945} remain the best short guide to the habits that the routine above tries to make deliberate.

\printchapterbib
\end{refsection}
